\cleardoublepage
\thispagestyle{empty}
\begin{center}


{\Large\textbf{Atom – Theorie und Anwendungen}}\\[1.2em]
{\large Dipl.-Ing.\,(FH)\,Christian Weilharter}\\[1.2em]
\textcopyright~2025, Christian Weilharter, Traunstein\\[2em]
\end{center}
\begin{flushleft}
	\begin{tabular}{@{}l l}
	
	
		\textbf{ISBN: (Print)} & 978-3-912302-14-1 \\[0.5em]
		%	\textbf{ISBN: (E-Book)} & 978-3-912302-01-1 \\[0.5em]
		\textbf{Auflage:} & 1.~Auflage 2026 \\[0.5em]
		\textbf{Satz:} & \LaTeX \\[0.5em]
		\textbf{Verlag:} & Christian Weilharter \\[0.5em]
		\textbf{Druck:} & Amazon KDP (Print-on-Demand) \\[0.5em]

		
		\textbf{Kontakt:} & \href{mailto:info@mathandphysics.de}{info@mathandphysics.de}\\[0.5em]
		\textbf{Web:} & \href{https://www.mathandphysics.de}{www.mathandphysics.de}\\
		

	\end{tabular}
\end{flushleft}

\vspace{2em}
\noindent
Alle Rechte vorbehalten. Kein Teil dieses Buches darf ohne schriftliche Genehmigung des Autors 
in irgendeiner Form reproduziert, gespeichert oder übertragen werden, 
weder elektronisch, mechanisch, durch Fotokopien, Aufnahmen noch auf andere Weise.
\begin{center}\small Printed in Germany\end{center}

\cleardoublepage
\chapter*{Vorwort}
\addcontentsline{toc}{chapter}{Vorwort}

Das Elektron gehört zu den unscheinbarsten und zugleich folgenreichsten Objekten der modernen Physik. 
Es ist unsichtbar klein, besitzt aber eine elektrische Ladung, eine Masse, einen Spin und eine zentrale Rolle beim Aufbau der Materie. 
Ohne Elektronen gäbe es keine Atome in ihrer heutigen Form, keine chemischen Bindungen, keine elektrische Leitung, keine Halbleiter, keine Computer, keine moderne Medizintechnik und keine Quantenelektrodynamik.


\begin{flushright}
	\textit{Dipl.-Ing.(FH) Christian Weilharter} \\
	\vspace{0.5em}
	[Traunstein, 2025]
\end{flushright}


